\section*{Hoofdstuk \ref{ch:b2:planet-formation-meteorites} --- De vorming van planeten en meteorieten}

\begin{solution}{exo:b2:planet-formation-meteorites:1}
Gas: \qty{2200}{g.cm^{-2}} op \qty{1}{au}; op \qty{5.2}{au}
$2200\times5.2^{-3/2} = 2200/11.86 = \qty{186}{g.cm^{-2}}$. Vaste stoffen:
$33\times0.33 = \qty{10.9}{g.cm^{-2}}$ op \qty{1}{au};
$33/11.86 = \qty{2.78}{g.cm^{-2}}$ op \qty{5.2}{au}. Op \qty{5.2}{au} is het gas
$186/2.78 = 67$ keer zwaarder dan de vaste stoffen.
\end{solution}

\begin{solution}{exo:b2:planet-formation-meteorites:2}
Aluminium (\qty{1653}{K}, ultrarefractair), magnesium (\qty{1336}{K}, gewoon),
ijzer (\qty{1334}{K}, gewoon), natrium (\qty{958}{K}, matig vluchtig), zwavel
(\qty{664}{K}, vluchtig), zuurstof als waterijs (\qty{180}{K}, sterk vluchtig).
CAI's, de eerste condensaten, zijn verrijkt in aluminium (en calcium, titaan);
de aarde is verarmd in de vluchtige elementen, natrium, zwavel en water, ten opzichte van
magnesium en aluminium.
\end{solution}

\begin{solution}{exo:b2:planet-formation-meteorites:3}
Vallen: gewone \omterm{def:b2:planet-formation-meteorites:chondrule}{chondrieten} $856/1107 = \qty{77.3}{\percent}$, \omterm{def:b2:planet-formation-meteorites:classes}{ijzermeteorieten}
$49/1107 = \qty{4.4}{\percent}$. Vondsten: $38647/44534 = \qty{86.8}{\percent}$ en
$1020/44534 = \qty{2.3}{\percent}$. Antarctische vondsten: $22096/44534 =
\qty{49.6}{\percent}$. De vallen: elke waargenomen val wordt verzameld, ongeacht zijn
klasse, terwijl een vondst afhangt van hoe goed de \omterm{def:b2:planet-formation-meteorites:fall-find}{meteoriet} op de grond bewaard blijft,
hoe hij eruitziet en waar men zoekt, en de helft van de vondsten komt van één soort
plaats.
\end{solution}

\begin{solution}{exo:b2:planet-formation-meteorites:4}
$\Delta t = 1.035\ln(\num{5.23e-5}/\num{1.0e-5}) = 1.035\ln 5.23 =
1.035\times1.654 = \qty{1.71}{Myr}$ na de CAI's.
\end{solution}

\begin{solution}{exo:b2:planet-formation-meteorites:5}
$4\pi\Sigma_1 r_1^2 = 4\pi\times\num{2.2e4}\times(\num{1.496e11})^2 =
\qty{6.19e27}{kg}$. Binnen \qty{30}{au}: $\times\sqrt{30} = 5.48$, dus
\qty{3.39e28}{kg} $= 0.017$ zonsmassa; binnen \qty{100}{au}: $\times10$,
\qty{6.19e28}{kg} $= 0.031$ zonsmassa. Beide liggen binnen het waargenomen bereik,
dichter bij de bovengrens: de nevel was een gewone schijf.
\end{solution}

\begin{solution}{exo:b2:planet-formation-meteorites:6}
$r = (120/T)^{7/3}$ au: voor water $(120/182)^{7/3} = 0.659^{7/3} =
\qty{0.38}{au}$; voor methaan $(120/41)^{7/3} = 2.93^{7/3} = \qty{12.3}{au}$.
De korrel in direct zonlicht is de oppervlaktelaag van de schijf; het middenvlak,
afgeschermd door het stof erboven, is kouder, dus ijs overleeft dichter bij de ster. De
\omterm{def:b2:planet-formation-meteorites:planetesimal}{planetesimalen} ontstonden in het middenvlak, dus de schatting voor het middenvlak geldt, op voorwaarde dat
de schijf passief is; verwarming door accretie schuift de \omterm{def:b2:planet-formation-meteorites:snow-line}{sneeuwlijn} naar buiten.
\end{solution}

\begin{solution}{exo:b2:planet-formation-meteorites:7}
Op \qty{1}{au}: $r^2 = \num{2.238e22}$, $\Sigma_p = \qty{109}{kg.m^{-2}}$, dus
$2\pi b r^2\Sigma_p = 31.42\times\num{2.238e22}\times109 = \num{7.66e25}$;
$(\num{7.66e25})^{3/2} = \num{6.71e38}$, gedeeld door
$(3\times\num{1.989e30})^{1/2} = \num{2.443e15}$: \qty{2.75e23}{kg}
$= 0.046$ aardmassa. Op \qty{5.2}{au}: $r^2 = \num{6.05e23}$,
$\Sigma_p = \qty{27.8}{kg.m^{-2}}$, $2\pi b r^2\Sigma_p = \num{5.29e26}$,
$M = \num{1.22e40}/\num{2.443e15} = \qty{4.98e24}{kg} = 0.83$ aardmassa.
Met \num{3.5} keer meer vaste stoffen groeit $M$ als $3.5^{3/2} = 6.55$: 5.5
aardmassa's.
\end{solution}

\begin{solution}{exo:b2:planet-formation-meteorites:8}
Per kilogram zijn er $\num{1.01e19}$ atomen, dus
$\num{1.01e19}\times\num{5.00e-13} = \qty{5.05e6}{J.kg^{-1}}$. De
temperatuurstijging zou $\num{5.05e6}/837 = \qty{6030}{K}$ zijn. Opwarmen van \num{250} tot
\qty{1700}{K} vergt $837\times1450 = \qty{1.21e6}{J.kg^{-1}}$: aluminium-26
leverde ongeveer vier keer zoveel.
\end{solution}

\begin{solution}{exo:b2:planet-formation-meteorites:9}
$\num{5.07e10}/\num{7.49e7} = 677$. De gemiddelde levensduur van ijzer-60 is
$2.62/\ln2 = \qty{3.78}{Myr}$. Aanvankelijke vermogens per gram:
$\num{5.07e10}/1.035 = \num{4.90e10}$ en $\num{7.49e7}/3.78 = \num{1.98e7}$ erg per
gram per miljoen jaar, een verhouding van ongeveer \num{2470}: aluminium-26 overheerste ruimschoots.
\end{solution}

\begin{solution}{exo:b2:planet-formation-meteorites:10}
De vermogens zijn gelijk wanneer $\num{4.90e10}\,e^{-t/1.035} = \num{1.98e7}\,e^{-t/3.78}$,
dat wil zeggen $t(1/1.035 - 1/3.78) = \ln(\num{4.90e10}/\num{1.98e7}) = \ln 2475 = 7.81$;
met $0.966 - 0.265 = 0.702$ per miljoen jaar is $t = \qty{11.1}{Myr}$. Tegen dan
heeft ijzer-60 nog $\num{7.49e7}\,e^{-11.1/3.78} = \num{7.49e7}\times0.053 =
\qty{4.0e6}{erg.g^{-1}}$ over, dat is \qty{400}{J.kg^{-1}}: duizenden keren te
weinig om wat dan ook te doen smelten.
\end{solution}

\begin{solution}{exo:b2:planet-formation-meteorites:11}
$\tau = 1.035\times\num{3.156e13} = \qty{3.267e13}{s}$ en
$H_0\tau = \qty{6.21e6}{J.kg^{-1}}$. Voor \qty{1200}{K}:
$t_a^* = 1.035\ln\bigl(\num{6.21e6}/(837\times950)\bigr) = 1.035\ln7.81 =
\qty{2.13}{Myr}$; voor \qty{1700}{K}: $1.035\ln\bigl(\num{6.21e6}/(837\times1450)\bigr)
= 1.035\ln5.12 = \qty{1.69}{Myr}$. Een halvering van $H_0$ trekt $\tau\ln2 =
\qty{0.72}{Myr}$ van beide af, want $t_a^* = \tau\ln H_0 + \tau\ln(\tau/c(T_c -
T_0))$ en alleen de eerste term verandert: de verschuiving is één halveringstijd, wat de
kritische temperatuur ook is.
\end{solution}

\begin{solution}{exo:b2:planet-formation-meteorites:12}
$\kappa = 2.1/(3300\times837) = \qty{7.60e-7}{m^2.s^{-1}}$ en
$\sqrt{\kappa\tau} = \sqrt{\num{7.60e-7}\times\num{3.267e13}} = \qty{4980}{m}$,
ongeveer \qty{5}{km}. Het centrum van een lichaam van \qty{100}{km} groot ligt \qty{50}{km} diep,
tien keer deze afstand: zijn warmte kan niet ontsnappen terwijl aluminium-26 vervalt, en de
propositie is van toepassing. Een lichaam van \qty{2}{km} verliest zijn warmte even snel als het
die maakt. Smelten vergt een straal van enkele keren \qty{5}{km}, dat wil zeggen een diameter van zo'n
\qtyrange{10}{20}{km} of meer, zoals de modellen vinden.
\end{solution}

\begin{solution}{pb:b2:planet-formation-meteorites:1}
\textbf{1.} Een val: zijn aankomst werd gezien en hij werd de volgende ochtend opgeraapt.
Hij is onverweerd, en hij vult de onvertekende steekproef aan van wat de aarde bereikt.

\textbf{2.} $856/1107 = 0.77$: ongeveer drie kansen op vier.

\textbf{3.} \omterm{def:b2:planet-formation-meteorites:chondrule}{Chondrulen} van \qty{0.3}{mm}, \qty{9}{\percent} metaal en olivijn met
\qty{18}{\percent} fayaliet: een \omterm{def:b2:planet-formation-meteorites:chondrule}{H-chondriet} (\qty{10}{\percent} metaal, Fa 16--20).

\textbf{4.} Vage \omterm{def:b2:planet-formation-meteorites:chondrule}{chondrulen} en uniforme mineralen wijzen op thermische metamorfose: types
5 tot 6. De gewone \omterm{def:b2:planet-formation-meteorites:chondrule}{chondrieten} bereikten hoogstens ongeveer \qty{1220}{K}, dat is
\qty{947}{\degreeCelsius}.

\textbf{5.} Gewone \omterm{def:b2:planet-formation-meteorites:chondrule}{chondrieten} bevatten slechts \qtyrange{0.1}{1}{\percent} CAI's in
volume. De steen bestaat uit \omterm{def:b2:planet-formation-meteorites:chondrule}{chondrulen} en matrix die na de CAI's zijn samengevoegd, en die
zijn de oudste vaste stoffen: elk gesteente dat ze bevat, is jonger.

\textbf{6.} Elk mineraal kreeg bij zijn vorming hetzelfde magnesium plus aluminium-26
in verhouding tot zijn aluminium; het aluminium-26 is sindsdien in
magnesium-26 overgegaan, dus hoe meer aluminium een mineraal per magnesium had, hoe meer
magnesium-26 het won.

\textbf{7.} $(0.139590 - 0.139500)/(15.0 - 0.05) = \num{0.000090}/14.95 =
\num{6.02e-6}$: de verhouding $\ce{^{26}Al}/\ce{^{27}Al}$ van de \omterm{def:b2:planet-formation-meteorites:chondrule}{chondrule} bij haar
vorming.

\textbf{8.} $0.139500 + \num{6.02e-6}\times3.00 = 0.139518$: op de rechte.

\textbf{9.} $\Delta t = 1.035\ln(\num{5.23e-5}/\num{6.02e-6}) = 1.035\ln8.69 =
1.035\times2.162 = \qty{2.24}{Myr}$.

\textbf{10.} $4567.5 - 2.2 = \qty{4565.3}{\Ma}$.

\textbf{11.} $\num{5.23e-5}\,e^{-10/1.035} = \num{5.23e-5}\times\num{6.4e-5} =
\num{3.3e-9}$, veel te weinig magnesium-26 om te meten. Hafnium-182 (halveringstijd
\qty{8.9}{Myr}) of mangaan-53 (\qty{3.7}{Myr}) is geschikt voor gebeurtenissen zo'n tien miljoen jaar
na de CAI's.

\textbf{12.} $\kappa = 2.1/(3300\times837) = \qty{7.60e-7}{m^2.s^{-1}}$;
$R^2/\kappa = \num{4e8}/\num{7.60e-7} = \qty{5.26e14}{s} = \qty{16.7}{Myr}$, zestien
keer de gemiddelde levensduur van aluminium-26: het centrum houdt zijn warmte vast.

\textbf{13.} $\num{1.9e-7}\,e^{-2.4/1.035} = \num{1.9e-7}\times0.0984 =
\qty{1.87e-8}{W.kg^{-1}}$.

\textbf{14.} $\num{1.87e-8}\times\num{3.267e13} = \qty{6.11e5}{J.kg^{-1}}$.

\textbf{15.} $\Delta T = \num{6.11e5}/837 = \qty{730}{K}$; piek
$250 + 730 = \qty{980}{K}$.

\textbf{16.} Boven \qty{750}{K}, dus het gesteente sinterde en werd gemetamorfoseerd; onder
\qty{1200}{K}, dus noch het metaal noch het sulfide smolt. De steen is een verhitte \omterm{def:b2:planet-formation-meteorites:chondrule}{chondriet},
van een hoog \omterm{def:b2:planet-formation-meteorites:type}{petrologisch type}, onder de \qty{1220}{K} van vraag 4: dat strookt.

\textbf{17.} $1.035\ln\bigl(\num{6.21e6}/(837\times950)\bigr) = \qty{2.13}{Myr}$:
ongeveer \qty{0.3}{Myr} vroeger.

\textbf{18.} De stijging is $\Delta T(1 - e^{-(t - t_a)/\tau})$; \qty{90}{\percent}
vergt $e^{-(t - t_a)/\tau} = 0.1$, $t - t_a = 1.035\ln10 = \qty{2.38}{Myr}$, dat is
\qty{4.8}{Myr} na de CAI's.

\textbf{19.} $\Delta T = 7420\,e^{-0.8/1.035} = 7420\times0.462 = \qty{3430}{K}$,
piek \qty{3680}{K}, ver boven de \qty{1700}{K} waarbij de silicaten gesmolten zijn.
Het lichaam smolt (de smeltwarmte hield zijn temperatuur laag) en
differentieerde in een metalen kern en een silicaatmantel en -korst.

\textbf{20.} $e^{(2.4 - 0.8)/1.035} = e^{1.546} = 4.7$ keer meer.

\textbf{21.} Zink en lood condenseren rond \qty{727}{K}, matig vluchtig; magnesium
(\qty{1336}{K}) en aluminium (\qty{1653}{K}) condenseren veel eerder. De vaste stoffen van
het gebied van de steen werden verzameld voordat het gas genoeg was afgekoeld om al het zink en
lood te laten condenseren.

\textbf{22.} $33\times0.33\times2.5^{-3/2} = 10.89/3.95 = \qty{2.76}{g.cm^{-2}}$,
\qty{27.6}{kg.m^{-2}}.

\textbf{23.} $r^2 = (3.74\times10^{11})^2 = \num{1.40e23}$;
$2\pi\times5\times\num{1.40e23}\times27.6 = \num{1.21e26}$;
$(\num{1.21e26})^{3/2}/\num{2.443e15} = \qty{5.45e23}{kg}$, $0.091$ aardmassa.
Het moederlichaam, $\tfrac43\pi(\num{2e4})^3\times3300 = \qty{1.1e17}{kg}$, is vijf
miljoen keer lichter: een \omterm{def:b2:planet-formation-meteorites:planetesimal}{planetesimaal}, geen embryo.

\textbf{24.} Het centrum van het moederlichaam piekte op ongeveer
\textbf{\qty[text-series-to-math=true]{707}{\degreeCelsius}} ($980 - 273$): heet genoeg om
de \omterm{def:b2:planet-formation-meteorites:chondrule}{chondriet} te laten rekristalliseren, te koel om zijn metaal te doen smelten.
\end{solution}
